\chapter{Curves}
\label{pa-ch-curves}

Let \(k\) be a field and \(C\to\Spec k\) a smooth, proper, geometrically
integral curve.

\begin{definition}[Genus]
  \label{pa-def-genus}
  \dcref{IV.1,custom}
  \cite{milne-av,kleiman-picard}
  The genus of \(C\) is
  \[
    g(C)=\dim_k H^1(C,\mathcal O_C).
  \]
\end{definition}

Proper coherent-cohomology finiteness makes the displayed dimension finite;
geometric integrality identifies it with the usual genus after any field
extension.

\begin{definition}[Divisors and degree]
  \label{pa-def-divisors}
  \dcref{II.6,custom}
  \cite{hartshorne-ag,papaioannou-rr}
  A divisor is \(D=\sum_xn_x[x]\) with finite support on the closed points of
  \(C\), and
  \[
    \deg D=\sum_xn_x[\kappa(x):k].
  \]
  Principal divisors are \(\operatorname{div}(f)\) for
  \(f\in K(C)^\times\); divisor classes form \(\operatorname{Pic}(C)\).
\end{definition}

\begin{theorem}[Riemann--Roch]
  \label{pa-thm-riemann-roch}
  \dcref{IV.1,custom}
  \uses{pa-def-genus,pa-def-divisors}
  \cite{hartshorne-ag,papaioannou-rr}
  There is a canonical divisor \(K_C\) of degree \(2g(C)-2\) and
  \[
    \dim_kH^0(C,\mathcal O_C(D))-
    \dim_kH^0(C,\mathcal O_C(K_C-D))
    =\deg D+1-g(C).
  \]
\end{theorem}

\begin{proof}
  Serre duality gives
  \(H^1(C,\mathcal O_C(D))^\vee\simeq
  H^0(C,\mathcal O_C(K_C-D))\).  For a closed point \(P\), the exact sequence
  \[
    0\longrightarrow\mathcal O_C(D-P)\longrightarrow\mathcal O_C(D)
    \longrightarrow\kappa(P)\longrightarrow0
  \]
  changes the Euler characteristic by \([\kappa(P):k]\).  Serre vanishing
  starts the induction at large degree, giving
  \(\chi(\mathcal O_C(D))=\deg D+\chi(\mathcal O_C)\).  Since
  \(\chi(\mathcal O_C)=1-g(C)\), duality yields the formula.
\end{proof}
